\chapter{2021年北京卷}

\section{单选题}

\begin{problem}
已知集合 \(A=\left\{x\mid -1<x<1\right\}\),\(B=\left\{x\mid 0\le x\le2\right\}\),则 \(A\cup B=\)(\quad).

\choices
    {\(\left\{x\mid -1<x<2\right\}\)}
    {\(\left\{x\mid -1<x\le2\right\}\)}
    {\(\left\{x\mid 0\le x<1\right\}\)}
    {\(\left\{x\mid 0\le x\le2\right\}\)}
\end{problem}

\begin{problem}
若复数 \(z\) 满足 \((1-\i)\cdot z=2\),则 \(z=\)(\quad).

\choices
    {\(-1-\i\)}
    {\(-1+\i\)}
    {\(1-\i\)}
    {\(1+\i\)}
\end{problem}

\begin{problem}
设函数 \(f(x)\) 的定义域为 \([0,1]\),则``\(f(x)\) 在区间 \([0,1]\) 上单调递增''是``\(f(x)\) 在区间 \([0,1]\) 上的最大值为 \(f(1)\)''的(\quad).

\choices
    {充分而不必要条件}
    {必要而不充分条件}
    {充分必要条件}
    {既不充分也不必要条件}
\end{problem}

\begin{problem}
某四面体的三视图如图所示,该四面体的表面积为(\quad).

\bitmapfigure[width=5.0cm]{img/2021/beijing/q4_fig1.png}

\choices
    {\(\dfrac{3}{2}+\dfrac{\sqrt3}{2}\)}
    {\(3+\sqrt3\)}
    {\(\dfrac{3}{2}+\sqrt3\)}
    {\(3+\dfrac{\sqrt3}{2}\)}
\end{problem}

\begin{problem}
若双曲线 \(\dfrac{x^2}{a^2}-\dfrac{y^2}{b^2}=1\) 的离心率为 \(2\),且过点 \((\sqrt2,\sqrt3)\),则双曲线的方程为(\quad).

\choices
    {\(2x^2-y^2=1\)}
    {\(x^2-\dfrac{y^2}{3}=1\)}
    {\(5x^2-3y^2=1\)}
    {\(\dfrac{x^2}{2}-\dfrac{y^2}{6}=1\)}
\end{problem}

\begin{problem}
《中国共产党党旗党徽制作和使用的若干规定》指出,中国共产党党旗为旗面缀有金黄色党徽图案的红旗,通用规格有五种. 这五种规格党旗的长 \(a_1\),\(a_2\),\(a_3\),\(a_4\),\(a_5\)(单位:\(\mathrm{cm}\))成等差数列,对应的宽为 \(b_1\),\(b_2\),\(b_3\),\(b_4\),\(b_5\)(单位:\(\mathrm{cm}\)),且长与宽之比都相等. 已知 \(a_1=288\),\(a_5=96\),\(b_1=192\),则 \(b_3=\)(\quad).

\choices
    {64}
    {96}
    {128}
    {160}
\end{problem}

\begin{problem}
函数 \(f(x)=\cos x-\cos2x\) 是(\quad).

\choices
    {奇函数,且最大值为 \(2\)}
    {偶函数,且最大值为 \(2\)}
    {奇函数,且最大值为 \(\dfrac{9}{8}\)}
    {偶函数,且最大值为 \(\dfrac{9}{8}\)}
\end{problem}

\begin{problem}
某一时段内,从天空降落到地面上的雨水,未经蒸发,渗漏,流失而在水面上积聚的深度,称为这个时段的降雨量(单位:\(\mathrm{mm}\)). \(24\,\mathrm{h}\) 降雨量的等级划分如下:

\begin{center}
\begin{tabular}{c|c}
等级 & \(24\,\mathrm{h}\) 降雨量(精确到 \(0.1\)) \\
\hline
\(\cdots\) & \(\cdots\) \\
小雨 & \(0.1\sim9.9\) \\
中雨 & \(10.0\sim24.9\) \\
大雨 & \(25.0\sim49.9\) \\
暴雨 & \(50.0\sim99.9\) \\
\(\cdots\) & \(\cdots\)
\end{tabular}
\end{center}

在综合实践活动中,某小组自制了一个底面直径为 \(200\,\mathrm{mm}\),高为 \(300\,\mathrm{mm}\) 的圆锥形雨量器. 若一次降雨过程中,该雨量器收集的 \(24\,\mathrm{h}\) 的雨水高度是 \(150\,\mathrm{mm}\)(如图所示),则这 \(24\,\mathrm{h}\) 降雨量的等级是(\quad).

\begin{center}
\bitmapfigure[width=4.9cm]{img/2021/beijing/q8_fig1.png}
\end{center}

\choices
    {小雨}
    {中雨}
    {大雨}
    {暴雨}
\end{problem}

\begin{problem}
已知直线 \(y=kx+m\)(\(m\) 为常数)与圆 \(x^2+y^2=4\) 交于点 \(M\),\(N\). 当 \(k\) 变化时,若 \(|MN|\) 的最小值为 \(2\),则 \(m=\)(\quad).

\choices
    {\(\pm1\)}
    {\(\pm\sqrt2\)}
    {\(\pm\sqrt3\)}
    {\(\pm2\)}
\end{problem}

\begin{problem}
已知 \(\{a_n\}\) 是各项均为整数的递增数列,且 \(a_1\ge3\). 若 \(a_1+a_2+a_3+\cdots+a_n=100\),则 \(n\) 的最大值为(\quad).

\choices
    {9}
    {10}
    {11}
    {12}
\end{problem}

\section{填空题}

\begin{problem}
在 \(\left(x^3-\dfrac{1}{x}\right)^4\) 的展开式中,常数项为\fillinblank{}.
\end{problem}

\begin{problem}
已知抛物线 \(y^2=4x\) 的焦点为 \(F\),点 \(M\) 在抛物线上,\(MN\) 垂直 \(x\) 轴于点 \(N\). 若 \(|MF|=6\),则点 \(M\) 的横坐标为\fillinblank{};\(\triangle MNF\) 的面积为\fillinblank{}.
\end{problem}

\begin{problem}
已知向量 \(\bs{a}\),\(\bs{b}\),\(\bs{c}\) 在正方形网格中的位置如图所示,若网格纸上小正方形的边长为 \(1\),则 \((\bs{a}+\bs{b})\cdot\bs{c}=\)\fillinblank{};\(\bs{a}\cdot\bs{b}=\)\fillinblank{}.

\FigureLayoutDeclare{img/2021/beijing/q13_fig1.png}{5.0cm}{34bp 38bp 39bp 32bp}
\bitmapfigure[width=0.46\linewidth]{img/2021/beijing/q13_fig1.png}
\end{problem}

\begin{problem}
若点 \(A(\cos\theta,\sin\theta)\) 关于 \(y\) 轴的对称点为 \(B\left(\cos\left(\theta+\dfrac{\pi}{6}\right),\sin\left(\theta+\dfrac{\pi}{6}\right)\right)\),则 \(\theta\) 的一个取值为\fillinblank{}.
\end{problem}

\begin{problem}
已知函数 \(f(x)=|\lg x|-kx-2\). 给出下列四个结论:

\begin{circlelist}
\item 当 \(k=0\) 时,\(f(x)\) 恰有 \(2\) 个零点;

\item 存在负数 \(k\),使得 \(f(x)\) 恰有 \(1\) 个零点;

\item 存在负数 \(k\),使得 \(f(x)\) 恰有 \(3\) 个零点;

\item 存在正数 \(k\),使得 \(f(x)\) 恰有 \(3\) 个零点.
\end{circlelist}

其中所有正确结论的序号是\fillinblank{}.
\end{problem}

\section{解答题}

\begin{problem}
在 \(\triangle ABC\) 中,\(c=2b\cos B\),\(C=\dfrac{2\pi}{3}\).

\begin{enumerate}
    \item 求 \(B\) 的大小;
    \item 再从下列三个条件中选择一个作为已知,使 \(\triangle ABC\) 存在且唯一确定,求 \(BC\) 边上中线的长.
\end{enumerate}

\begin{circlelist}
\item \(c=\sqrt2b\);

\item \(\triangle ABC\) 的周长为 \(4+2\sqrt3\);

\item \(\triangle ABC\) 的面积为 \(\dfrac{3\sqrt3}{4}\).
\end{circlelist}
\end{problem}

\begin{problem}
如图,在正方体 \(ABCD-A_1B_1C_1D_1\) 中,\(E\) 为 \(A_1D_1\) 的中点,\(B_1C_1\) 与平面 \(CDE\) 交于点 \(F\).

\begin{enumerate}
    \item 求证:\(F\) 为 \(B_1C_1\) 中点;
    \item 若 \(M\) 是棱 \(A_1B_1\) 上一点,且二面角 \(M-FC-E\) 的余弦值为 \(\dfrac{\sqrt5}{3}\),求 \(\dfrac{A_1M}{A_1B_1}\) 的值.
\end{enumerate}

\bitmapfigure[width=6.0cm]{img/2021/beijing/q17_fig1.png}
\end{problem}

\begin{problem}
在核酸检测中,``\(k\) 合 1''混采核酸检测是指:先将 \(k\) 个人的样本混合在一起进行 1 次检测,如果这 \(k\) 个人都没感染新冠病毒,则检测结果为阴性;如果这 \(k\) 个人中有感染新冠病毒,则检测结果为阳性,此时需对每人再进行一次检测,得到每人的检测结果,检测结束. 现对 \(100\) 人进行核酸检测,假设其中只有 \(2\) 人感染新冠病毒,并假设每次检测结果准确.

\begin{enumerate}
    \item 将这 \(100\) 人随机分成 \(10\) 组,每组 \(10\) 人,且对每组都采用``10 合 1''混采核酸检测.
    \begin{enumerate}
        \item 如果感染新冠病毒的 \(2\) 人在同一组,求检测的总次数;
        \item 已知感染新冠病毒的 \(2\) 人在同一组的概率为 \(\dfrac{1}{11}\),设 \(X\) 是检测的总次数,求 \(X\) 的分布列与数学期望 \(EX\);
    \end{enumerate}
    \item 将这 \(100\) 人随机分成 \(20\) 组,每组 \(5\) 人,且对每组都采用``5 合 1''混采核酸检测. 设 \(Y\) 是检测的总次数,试判断数学期望 \(EY\) 与(1)中的 \(EX\) 的大小.(结论不要求证明)
\end{enumerate}
\end{problem}

\begin{problem}
已知函数 \(f(x)=\dfrac{3-2x}{x^2+a}\).

\begin{enumerate}
    \item 若 \(a=0\),求曲线 \(y=f(x)\) 在点 \((1,f(1))\) 处的切线方程;
    \item 若 \(f(x)\) 在 \(x=-1\) 处取得极值,求 \(f(x)\) 的单调区间,并求其最大值与最小值.
\end{enumerate}
\end{problem}

\begin{problem}
已知椭圆 \(E:\dfrac{x^2}{a^2}+\dfrac{y^2}{b^2}=1\)(\(a>b>0\))的一个顶点为 \(A(0,-2)\),以椭圆 \(E\) 的四个顶点为顶点的四边形面积为 \(4\sqrt5\).

\begin{enumerate}
    \item 求椭圆 \(E\) 的方程;
    \item 过点 \(P(0,-3)\) 斜率为 \(k\) 的直线 \(l\) 与椭圆 \(E\) 交于不同的两点 \(B\),\(C\),直线 \(AB\),\(AC\) 分别与直线 \(y=-3\) 交于 \(M\),\(N\). 当 \(|PM|+|PN|\le15\) 时,求 \(k\) 的取值范围.
\end{enumerate}
\end{problem}

\begin{problem}
设 \(p\) 为实数. 若无穷数列 \(\{a_n\}\) 满足如下三个性质,则称 \(\{a_n\}\) 为 \(R_p\) 数列:

\begin{circlelist}
\item \(a_1+p\ge0\),且 \(a_2+p=0\);

\item \(a_{4n-1}<a_{4n}\)(\(n=1\),\(2\),\(\ldots\));

\item \(a_{m+n}\in\left\{a_m+a_n+p,\ a_m+a_n+p+1\right\}\)(\(m=1\),\(2\),\(\ldots\);\(n=1\),\(2\),\(\ldots\)).
\end{circlelist}

\begin{enumerate}
    \item 如果数列 \(\{a_n\}\) 的前四项分别为 \(2\),\(-2\),\(-2\),\(-1\),那么 \(\{a_n\}\) 是否可能为 \(R_2\) 数列?说明理由;
    \item 若数列 \(\{a_n\}\) 是 \(R_0\) 数列,求 \(a_5\);
    \item 设数列 \(\{a_n\}\) 的前 \(n\) 项和为 \(S_n\),是否存在 \(R_p\) 数列 \(\{a_n\}\),使得 \(S_n\ge S_{10}\) 恒成立?如果存在,求出所有的 \(p\);如果不存在,说明理由.
\end{enumerate}
\end{problem}

